\begin{tabular}{p{2.6cm}p{3.6cm}p{3.0cm}p{4.2cm}p{4.4cm}p{4.6cm}}
\toprule
\rowcolor{naranja}
\textbf{Estudio} & \textbf{(a) Ámbito y granularidad} & \textbf{(b) Horizonte temporal} & \textbf{(c) Matriz exógena} & \textbf{(d) Protocolo de validación} & \textbf{(e) Aporte empírico} \\
\midrule
\rowcolor{filaclara} Monje et al. (2022) & Una línea de autobús (EMT línea 1), solo franja de tarde & 26 meses pre-pandemia (ene. 2015 - feb. 2017) & Festivo binario + calendario (mes, día de semana); cero meteorología & Holdout cronológico simple 80/\allowbreak{}20, una sola ejecución, sin CV & Reporta el ganador (LSTM R²=0,89) sin MAE/\allowbreak{}RMSE/\allowbreak{}MAPE absolutos \\
Toqué et al. (2017) & Transporte multimodal por estación (París) & Logs de billetaje, 2017 & No especificada en la fuente disponible & No especificado en la fuente disponible & Comparativa de métodos de ML para flujos multimodales, sin arquitectura híbrida residual \\
\rowcolor{filaclara} Cardozo et al. (2012) & Entradas por estación del Metro de Madrid, corte transversal espacial & Sin dimensión temporal (regresión espacial) & Entorno construido; sin calendario ni clima & Ajuste espacial in-sample & Regresión espacial por estación; no aborda predicción temporal \\
Zhang (2003) & Series univariantes ajenas al transporte (manchas solares, linces, GBP/\allowbreak{}USD) & 288 /\allowbreak{} 114 /\allowbreak{} 731 observaciones & Ninguna (residuo modelado a partir de sus propios rezagos) & Holdout único, residuos ARIMA in-sample (sin OOF) & El híbrido ARIMA+red neuronal bate a ambos componentes en las tres series \\
\rowcolor{filaclara} Surribas-Sayago et al. [ficha sin verificar] & Radiación solar difusa (dominio meteorológico, no transporte) & No verificado & Retardos y términos de Fourier (fuente de inspiración de la ingeniería de variables de este TFM, no de su arquitectura) & No verificado & Arquitectura paralela CNN+LSTM+MLP; citada solo por su ingeniería de características \\
Este TFM & Demanda diaria agregada de los cuatro operadores del CRTM (Madrid) & 1.310 días continuos 2023-2026, cero huecos, cero nulos, íntegramente post-pandemia & 161 variables: 105 meteorológicas físicas retardadas (lags 1/\allowbreak{}2/\allowbreak{}3/\allowbreak{}7 + media móvil 7), 6 términos de Fourier, 55 días puente detectados; meteorología del mismo día explícitamente prohibida & Partición cronológica 70/\allowbreak{}15/\allowbreak{}15 de frontera única; walk-forward OOF en 5 bloques expansivos; SARIMAX walk-forward de un paso sin reestimación; escalado ajustado solo en train; 257 tests fijando las reglas anti-fuga & Auditoría del fallo: el híbrido no bate a SARIMAX (+70.597 MAE test) ni a XGBoost solo (+78.103); el mejor modelo es un ensamble 50/\allowbreak{}50 (MAE 139.725, -10,5\% vs. el mejor componente) \\
\bottomrule
\end{tabular}
